% TOP_PROBLEM: {"id": 30007243, "title": "A parity lower bound for bounded-depth quantum circuits with polynomial ancillas", "category_id": 16, "category": {"id": 16, "name": "physics", "display_name": "Mathematical Physics", "description": "Problems at the intersection of mathematics and physics.", "slug": "physics", "order_index": 16, "created_at": "2026-07-31T15:26:25.671Z"}, "status": "open"}
% FIRST_POSTED: 2026-10-06
% ATTEMPT_STATUS: solved
% OPENAI_CLAIM: {"schema_version":1,"source_repo":"https://github.com/openai/math","source_commit":"adc7f1241b42e322a6451854ab7e4b4c146bf78a","release_date":"2026-10-06","match":"full","resolution":"proved","independently_reviewed":false,"second_pass":"agreed","families":[{"family":"274","title":"Parity is not in QAC<sup>0</sup>","subject":"Mathematical physics","manuscripts":[{"title":"Product-projection localization and the QAC ⁰ parity lower bound","pdf_path":"preprints/Product-projection-localization-and-the-QAC0-parity-lower-bound-September-24-2026/paper.pdf","readme_path":"preprints/Product-projection-localization-and-the-QAC0-parity-lower-bound-September-24-2026/README.md","date":"2026-09-24","theorem_ref":"thm:parity (sections/introduction.tex:25)","statement_sources":[{"path":"preprints/Product-projection-localization-and-the-QAC0-parity-lower-bound-September-24-2026/build/sections/introduction.tex","line":25,"label":"thm:parity"}]},{"title":"Regular trajectories, pruning and quantum parity","pdf_path":"preprints/Regular-trajectories-pruning-and-quantum-parity-September-24-2026/paper.pdf","readme_path":"preprints/Regular-trajectories-pruning-and-quantum-parity-September-24-2026/README.md","date":"2026-09-24","theorem_ref":"thm:parity (sections/01-introduction.tex:47)","statement_sources":[{"path":"preprints/Regular-trajectories-pruning-and-quantum-parity-September-24-2026/build/sections/01-introduction.tex","line":47,"label":"thm:parity"}]}],"lean":{"doc_path":"lean/docs/274.md","comparators":["lean/ComparatorChallenges/QACParity.json","lean/ComparatorChallenges/QACParity.lean","lean/ComparatorChallenges/RegularParity.json","lean/ComparatorChallenges/RegularParity.lean"],"declarations":["OAI.QAC.parity_lower_bound","OAI.QAC.parity_lower_bound_polynomial_size"],"covers_main_theorem":true},"reasoning_trace":null}]}
\documentclass[11pt]{article}
\usepackage[margin=25mm]{geometry}
\usepackage{fontspec}
\setmainfont{Latin Modern Roman}
\usepackage{amsmath,amssymb,mathtools}
\usepackage{unicode-math}
\setmathfont{Latin Modern Math}
\usepackage{enumitem}
\usepackage{xurl,hyperref,bookmark}
\hypersetup{colorlinks=true,urlcolor=blue}
\setcounter{secnumdepth}{0}
\setlength{\parindent}{0pt}
\setlength{\parskip}{5pt}
\setlength{\emergencystretch}{3em}
\begin{document}
\section*{A parity lower bound for bounded-depth quantum circuits with polynomial ancillas}

\subsection{OpenAI's claim}
The linked manuscripts present the following claims by OpenAI: Polynomial-qubit constant-depth circuits of one-qubit and unbounded Toffoli gates are claimed unable to compute parity with fixed worst-case advantage; A separate regular-trajectory proof claims the same measured-output parity lower bound for $\mathrm{QAC}^0$ circuits with zero-initialized ancillas. Together they give a claimed affirmative answer in the matching scope.

\subsection{Mathematical statement of the claimed result}
\paragraph{Result 1.} Fix depth bound \(d\ge0\), ancilla exponent \(c\ge1\), and advantage \(0<\varepsilon\le1/2\). Use \(n\) input qubits, with a total register size bounded by \(n^c\). Permit arbitrary unitaries acting on single qubits, together with Toffoli gates of any arity, with disjoint supports in each layer, zero-initialized ancillas, and one measured output. The claim is that for all sufficiently large \(n\), every circuit of depth at most \(d\) has an input \(x\) with
\[
 \mathbb P(\text{output}=x_1\oplus\cdots\oplus x_n)<1/2+\varepsilon.
\]
Unused output registers may be discarded.

\paragraph{Result 2.} Fix positive integers \(d,k,C\) and \(0<\varepsilon\le1/2\). The allowed circuits use arbitrary one-qubit unitaries, unbounded Toffoli gates, disjoint supports within layers, zero ancillas, and a measured output bit. The source claims a finite \(n_0\) such that for every \(n\ge n_0\), any circuit with a total register size bounded by \(C(n+1)^k\), and with depth bounded by \(d\), has some input \(x\) satisfying
\[
 \mathbb P(\text{answer}=\operatorname{Parity}_n(x))<1/2+\varepsilon.
\]

\subsection{Outline of the argument}
Product projections localize the output test onto small supports. A controlled-coefficient polynomial approximation permits depth reduction, and the resulting low-degree description has insufficient parity correlation even when final garbage registers are unrestricted.

Two polynomial filters regularize circuit trajectories. Low-degree perturbations and excitation tails are transported through the gates, then projection insertion and pruning restrict the histories that can affect the output; this restriction forces the parity advantage to vanish.

\subsection{Scope relative to this problem}
Polynomial-qubit constant-depth circuits of one-qubit and unbounded Toffoli gates are claimed unable to compute parity with fixed worst-case advantage. The question retains these same parameters and conventions. A separate regular-trajectory proof claims the same measured-output parity lower bound for $\mathrm{QAC}^0$ circuits with zero-initialized ancillas. The question retains these same parameters and conventions. Within these assumptions, this is a claimed affirmative answer to the stated question.

\subsection{Formal verification}
For family 274, the documented entries include the measured-output parity lower bound for polynomial-qubit constant-depth circuits and the supporting product-projection estimate, with zero ancillas and unrestricted discarded registers. The listed Comparator files are \texttt{QACParity.lean}, \texttt{RegularParity.lean}. Named declarations include \texttt{OAI.QAC.parity\_lower\_bound}, \texttt{OAI.QAC.parity\_lower\_bound\_polynomial\_size}. This site has not executed the Lean code. The scope descriptions do not certify the remaining manuscript arguments.

\subsection{Status caveat}
OpenAI attributes these manuscripts to work by an internal model. Its release cautions that arguments without formalization can contain mistakes. This catalogue has not independently checked the proofs; the labels report the mathematical claims attributed to OpenAI.

\subsection{Sources}
\begin{itemize}
\item OpenAI, \emph{Product-projection localization and the QAC \({}^{0}\) parity lower bound}, 2026-09-24; citation key \texttt{OAI:Product-projection-localization-and-the-QAC0-parity-lower-bound-September-24-2026}. \href{https://github.com/openai/math/blob/adc7f1241b42e322a6451854ab7e4b4c146bf78a/preprints/Product-projection-localization-and-the-QAC0-parity-lower-bound-September-24-2026/paper.pdf}{PDF (GitHub)}; \href{https://github.com/openai/math/raw/adc7f1241b42e322a6451854ab7e4b4c146bf78a/preprints/Product-projection-localization-and-the-QAC0-parity-lower-bound-September-24-2026/paper.pdf}{PDF (direct)}; \href{https://github.com/openai/math/blob/adc7f1241b42e322a6451854ab7e4b4c146bf78a/preprints/Product-projection-localization-and-the-QAC0-parity-lower-bound-September-24-2026/README.md}{Manuscript page and citation}; \href{https://github.com/openai/math/blob/main/preprints/Product-projection-localization-and-the-QAC0-parity-lower-bound-September-24-2026/paper.pdf}{Latest version}; Statement source: \href{https://github.com/openai/math/blob/adc7f1241b42e322a6451854ab7e4b4c146bf78a/preprints/Product-projection-localization-and-the-QAC0-parity-lower-bound-September-24-2026/build/sections/introduction.tex\#L25}{thm:parity}.
\item OpenAI, \emph{Regular trajectories, pruning and quantum parity}, 2026-09-24; citation key \texttt{OAI:Regular-trajectories-pruning-and-quantum-parity-September-24-2026}. \href{https://github.com/openai/math/blob/adc7f1241b42e322a6451854ab7e4b4c146bf78a/preprints/Regular-trajectories-pruning-and-quantum-parity-September-24-2026/paper.pdf}{PDF (GitHub)}; \href{https://github.com/openai/math/raw/adc7f1241b42e322a6451854ab7e4b4c146bf78a/preprints/Regular-trajectories-pruning-and-quantum-parity-September-24-2026/paper.pdf}{PDF (direct)}; \href{https://github.com/openai/math/blob/adc7f1241b42e322a6451854ab7e4b4c146bf78a/preprints/Regular-trajectories-pruning-and-quantum-parity-September-24-2026/README.md}{Manuscript page and citation}; \href{https://github.com/openai/math/blob/main/preprints/Regular-trajectories-pruning-and-quantum-parity-September-24-2026/paper.pdf}{Latest version}; Statement source: \href{https://github.com/openai/math/blob/adc7f1241b42e322a6451854ab7e4b4c146bf78a/preprints/Regular-trajectories-pruning-and-quantum-parity-September-24-2026/build/sections/01-introduction.tex\#L47}{thm:parity}.
\item Family 274: \href{https://github.com/openai/math/blob/adc7f1241b42e322a6451854ab7e4b4c146bf78a/lean/docs/274.md}{Lean scope}.
\item \href{https://github.com/openai/math/blob/adc7f1241b42e322a6451854ab7e4b4c146bf78a/lean/ComparatorChallenges/QACParity.json}{Comparator QACParity.json}.
\item \href{https://github.com/openai/math/blob/adc7f1241b42e322a6451854ab7e4b4c146bf78a/lean/ComparatorChallenges/QACParity.lean}{Comparator QACParity.lean}.
\item \href{https://github.com/openai/math/blob/adc7f1241b42e322a6451854ab7e4b4c146bf78a/lean/ComparatorChallenges/RegularParity.json}{Comparator RegularParity.json}.
\item \href{https://github.com/openai/math/blob/adc7f1241b42e322a6451854ab7e4b4c146bf78a/lean/ComparatorChallenges/RegularParity.lean}{Comparator RegularParity.lean}.
\item \href{https://github.com/openai/math/blob/adc7f1241b42e322a6451854ab7e4b4c146bf78a/CONTENTS.md}{OpenAI catalogue entry}.
\end{itemize}
\end{document}
